% TOP_PROBLEM: {"id": 2531, "title": "Kourovka Notebook Problem 21.22", "category_id": 17, "category": {"id": 17, "name": "group_theory", "display_name": "Group Theory", "description": "Problems about groups, group actions, representations, and related algebraic structures.", "slug": "group-theory", "order_index": 17, "created_at": "2026-07-31T15:26:25.671Z"}, "status": "open", "problem_number": "KOU-21.22", "set_id": 10}
% FIRST_POSTED: 2026-10-01
% ENTRY_KIND: statement_only
% UnsolvedMath ID 2531; source catalogue code KOU-21.22
% !TeX program = lualatex
% Definition and sources only; this file is not an LLM solution attempt.
\documentclass[11pt,a4paper]{article}
\usepackage[margin=25mm]{geometry}
\usepackage{fontspec}
\setmainfont{lmroman10-regular.otf}[BoldFont=lmroman10-bold.otf,
 ItalicFont=lmroman10-italic.otf,BoldItalicFont=lmroman10-bolditalic.otf]
\usepackage{amsmath,amssymb,mathtools}
\usepackage{unicode-math}
\setmathfont{latinmodern-math.otf}
\AtBeginDocument{\let\setminus\smallsetminus}
\usepackage{xurl,hyperref}
\hypersetup{colorlinks=true,urlcolor=blue}
\setcounter{secnumdepth}{0}
\setlength{\parindent}{0pt}
\setlength{\parskip}{5pt}
\setlength{\emergencystretch}{3em}
\begin{document}
\section*{Kourovka Notebook Problem 21.22}\label{2531}
{\small UnsolvedMath ID 2531 · KOU-21.22 · Group Theory · Kourovka Notebook - New Problems, Issue 21}
\subsection{Definitions and mathematical statement}
A group is Hopfian if every surjective endomorphism of it is injective,
hence an automorphism. It is finitely generated if some finite set generates
it using multiplication and inverses.
Let $G$ and $H$ be finitely generated Hopfian groups.

Define the restricted base group
\[
 B=\bigoplus_{h\in H}G
   =\{f:H\to G:\{h:f(h)\ne1\}\text{ is finite}\},
\]
with pointwise multiplication. Distinct copies of $G$ commute inside
this direct sum, even when $G$ itself is nonabelian. The group $H$ acts
on $B$ by left translation of the indices:
$(h\cdot f)(x)=f(h^{-1}x)$. The standard restricted wreath product is
$G\wr H=B\rtimes H$, with multiplication
\[
 (f,h)(g,k)=(f\,(h\cdot g),hk).
\]
The question is whether $G\wr H$ must be Hopfian for every such pair.
Thus every surjective homomorphism $\Phi:G\wr H\to G\wr H$ would have
to have trivial kernel.

The map $\Phi$ is arbitrary; it is not assumed to preserve the displayed
base group, the embedded acting group, or the projection onto $H$.
The restricted support condition is essential: replacing $B$ by all
functions $H\to G$ gives a different wreath product. No abelianity,
residual finiteness, torsion condition, or finite-presentation hypothesis
is added to the factors. A counterexample requires both factors to be
Hopfian even though their restricted wreath product is not.
\subsection{Short English statement}
Is the standard restricted wreath product of two finitely generated Hopfian groups necessarily Hopfian?
\subsection{Sources}
\begin{itemize}
\item[S1] ULAM.AI and the UnsolvedMath Contributors, supplied problems.json, record 2531 (KOU-21.22), Kourovka Notebook - New Problems, Issue 21. Catalogue identity and source transcription; CC BY 4.0.
\url{https://huggingface.co/datasets/ulamai/UnsolvedMath}.
\item[S2] The Kourovka Notebook, 21st edition, new problems, Problem 21.22. Expanded formulation of the supplied catalogue statement.
\url{https://arxiv.org/pdf/1401.0300v47}.
\end{itemize}
\end{document}
